\chapter{Instrumentación del frontend}
\label{cap:frontend-instr}

\begin{intuicion}
La promesa plug-and-play (capítulo anterior) dice ``cada
optimización se nota inmediatamente''. Para que ``se note'', el
frontend tiene que mostrarla. La instrumentación que GraphHunter
expone vive en tres canales: \texttt{cmd\_get\_perf\_status}
(qué optimizaciones están activas), \texttt{cmd\_get\_memory\_status}
(presión de memoria), y la respuesta de \texttt{run\_hunt\_batch}
(cache hits del último batch). Las pildoras del toolbar resumen
los tres en un vistazo. Este capítulo cierra el libro mostrando
cómo se materializa la promesa.
\end{intuicion}

\section{Las pildoras del toolbar}

\begin{figure}[H]
\centering
\begin{tikzpicture}[
    node distance=0.4cm,
    every node/.style={font=\scriptsize},
    pill/.style={draw, rounded corners=4mm, minimum width=2.0cm,
                 minimum height=0.7cm, font=\footnotesize}
]
\node[pill, draw=accent, fill=accentlight] (perf) {Perf: M+L1};
\node[pill, draw=green!50!black, fill=green!10, right=of perf]
    (memok) {Mem: 14\%};
\node[pill, draw=orange!70!black, fill=orange!10, right=of memok]
    (memwarn) {Mem: 78\%};
\node[pill, draw=red!70!black, fill=red!10, right=of memwarn]
    (memcrit) {Mem: 92\%};
\end{tikzpicture}
\caption{Pildoras del toolbar. La de \emph{perf} indica si los
flags M+L1 están activos (verde) o si el usuario forzó el legacy
path (gris). La de \emph{mem} cambia de verde a ámbar
(\(\geq 75\%\) RSS/total) a rojo (\(\geq 90\%\)).}
\label{fig:plug-pills}
\end{figure}

\section{\texttt{cmd\_get\_perf\_status}}

El comando Tauri que el frontend invoca al mount:

\begin{lstlisting}[language=Rust]
#[tauri::command]
fn cmd_get_perf_status() -> PerfStatus {
    PerfStatus {
        lftj_planner: env_bool("GRAPHHUNTER_LFTJ_PLANNER", true),
        k4_lftj_marker: env_bool("GRAPHHUNTER_K4_LFTJ", true),
        yannakakis: env_bool("GRAPHHUNTER_YANNAKAKIS", true),
        // build capabilities
        has_nauty: cfg!(feature = "symmetry-nauty"),
        has_simd: cfg!(feature = "simd-intersect"),
    }
}
\end{lstlisting}

\noindent
La struct se serializa a JSON y el frontend la lee:

\begin{lstlisting}[language=Rust]
type PerfStatus = {
    lftj_planner: boolean;
    k4_lftj_marker: boolean;
    yannakakis: boolean;
    has_nauty: boolean;
    has_simd: boolean;
};
\end{lstlisting}

\subsection{Render de la pildora}

\filepath{apps/tauri/src/components/PerfPill.tsx}:

\begin{itemize}
    \item Si todos los flags estructurales (\texttt{lftj\_planner},
    \texttt{yannakakis}) están \texttt{true}: pildora ``Perf: M+L1''
    en verde acento.
    \item Si alguno está \texttt{false}: pildora ``Perf: legacy''
    en gris.
    \item Tooltip on-hover: lista los flags individuales y su
    estado actual.
\end{itemize}

\section{\texttt{cmd\_get\_memory\_status}}

\begin{lstlisting}[language=Rust]
#[tauri::command]
fn cmd_get_memory_status() -> MemoryStatus {
    let process = sysinfo::Process::current();
    let host = sysinfo::System::new();
    MemoryStatus {
        rss_bytes: process.memory(),
        host_total_bytes: host.total_memory(),
        ratio: process.memory() as f64 / host.total_memory() as f64,
    }
}
\end{lstlisting}

\noindent
El frontend lo polea cada 5 segundos y actualiza la pildora:

\begin{itemize}
    \item Ratio \(< 0.5\): verde, sin warning.
    \item Ratio \(\in [0.75, 0.90)\): ámbar, ``mem alta''. Sugiere
    cerrar otras apps o reducir el dataset.
    \item Ratio \(\geq 0.90\): rojo, ``crítica''. Recomienda
    \texttt{session\_save} y restart.
\end{itemize}

\noindent
La idea: el analista ve la presión \emph{antes} de que el motor
entre en swap (que serían minutos de UI freeze).

\section{Telemetría adicional en \texttt{run\_hunt\_batch}}

El endpoint \texttt{run\_hunt\_batch} (L2) ---para correr varios
hunts en una pasada compartiendo bag/result caches--- retorna una
respuesta extendida con cache counters:

\begin{lstlisting}[language=Rust]
pub struct RunHuntBatchResponse {
    pub entries: Vec<RunHuntBatchEntry>,
    pub indexes_were_warm: bool,
    pub bag_cache_hits: usize,         // L2 stage 2
    pub bag_cache_misses: usize,
    pub result_cache_hits: usize,      // L2 stage 3a
    pub last_prefix_resume_depth: usize,  // L2 stage 3b
    pub prefix_cache_hits: usize,
}
\end{lstlisting}

\noindent
Cada campo es una métrica que el frontend muestra como tooltip
sobre la pildora de perf. El analista que pasa el mouse ve:
\begin{itemize}
    \item ``Bag cache: 87\% hit ratio''
    \item ``Result cache: 12 hits, 1 miss''
    \item ``Last prefix resumed at depth 2''
\end{itemize}

\noindent
Si los hits son altos, las optimizaciones L2 están trabajando. Si
son cero, hay algo que investigar (¿hipótesis tan distintas que no
comparten nada? ¿caches sub-dimensionados?).

\section{Doc autocompliance}

El procedimiento de agregar un flag nuevo:

\begin{enumerate}
    \item Implementar la optimización en \texttt{core}.
    \item Agregar default-on en \filepath{apps/tauri/src-tauri/src/lib.rs}
    y \filepath{apps/cli/src/main.rs}.
    \item Extender \texttt{PerfStatus} con el nuevo campo.
    \item Extender \texttt{PerfPill.tsx} para reflejar el flag.
    \item Actualizar \filepath{08-env-vars.tex} y
    \filepath{12-preguntas-frecuentes.tex}.
    \item \texttt{make all} para regenerar PDFs.
    \item Si pasa CI con los seis PDFs frescos $\Rightarrow$
    merge.
\end{enumerate}

\noindent
El último punto es self-enforcing: si la doc no se regeneró, CI
falla.

\section{El cierre del libro}

Esta es la última pieza pedagógica del diseño. Los próximos pasos
para el lector:

\begin{itemize}
    \item \emph{Manual} (libro paralelo): cómo el analista usa el
    motor en su workflow diario.
    \item \emph{Especificación}: contracts HTTP, MCP, FFI, formatos
    de log soportados.
    \item \emph{Arquitectura}: estructura de crates, ADRs,
    decisiones de \emph{repo organization}.
    \item \emph{Paper}: justificación académica con benchmarks y
    comparaciones contra el estado del arte.
\end{itemize}

\noindent
Este libro es la guía pedagógica que une los cinco. Si llegaste
hasta acá, deberías poder leer cualquiera de los otros sin
dificultad ---y, más importante, contribuir al motor con
optimizaciones nuevas siguiendo los patrones que el libro documenta.

\section{Resumen}

\begin{itemize}
    \item Tres comandos de instrumentación:
    \texttt{cmd\_get\_perf\_status},
    \texttt{cmd\_get\_memory\_status}, telemetría en
    \texttt{run\_hunt\_batch}.
    \item Pildoras del toolbar: perf (verde/gris), mem
    (verde/ámbar/rojo).
    \item Cache counters en batch: visible vía tooltip.
    \item Procedimiento ``agregar un flag nuevo'' obliga a
    actualizar PDFs; CI lo enforce.
\end{itemize}

\begin{ejercicios}
\begin{enumerate}
    \item \starred{} Diseñá una pildora más para el toolbar:
    ``Build: AVX2'' que indique si el binary tiene SIMD activado.
    ¿Qué necesita en \texttt{PerfStatus}?

    \item Si la pildora de mem se vuelve roja durante un hunt,
    ¿qué debería hacer el frontend? Discutí: ¿abortar
    automáticamente? ¿alertar y dejar al usuario decidir?

    \item \starred{} El polling de memoria es cada 5 segundos.
    ¿Por qué no más frecuente? ¿Por qué no menos?

    \item Diseñá una métrica \emph{nueva} que valga la pena
    exponer: por ejemplo, ``aristas/seg que está procesando el
    ingest''. ¿Dónde la calculás? ¿Cómo la mostrás?

    \item ¿Por qué la doc se regenera en CI y no manualmente
    cuando el dev quiere? Pensá: ¿qué se rompe sin regeneración
    automática?
\end{enumerate}
\end{ejercicios}
